% Generated by tools/textbook/export_listings.js from decks/plan.js and code/l05.S: do not edit.
\section{Il programma di questo capitolo}\label{sec:lab:l05}

I listati di questo capitolo sono parti di un solo programma,
\fname{code/l05.S} nel repository del corso, stampato per intero nel
Listato~\ref{lst:lab:l05}. I test automatici del materiale del corso lo
compilano e lo eseguono sull'esecutore senza interfaccia di \prgqt{}, con
ingressi programmati, e controllano il frame; la
Figura~\ref{fig:lab:l05} è il frame scritto dal test.

\begin{description}[style=nextline,leftmargin=1.5em]
\item[Sullo schermo.] I 32 bit di marker\_x come lampade. Parte da 2 (una lampada); ogni frame con LEFT premuto sottrae 1: dopo 2, 1 e 0 si accendono tutte le lampade, 0xFFFFFFFF, cioè -1.
\item[Che cosa premere.] Tenete premuto LEFT per qualche frame.
\end{description}

\begin{figure}[htb]
\centering
\includegraphics[width=0.5\textwidth]{lab_l05.png}
\caption{\fname{l05.S} su \prgqt{}.}\label{fig:lab:l05}
\end{figure}

\asmfile[caption={Il programma completo del capitolo, \fname{code/l05.S}.},label=lst:lab:l05]{code/l05/program.S}

Per eseguirlo su \prgqt{} i comandi sono quelli del
Listato~\ref{lst:lab:l05:run}; il Paragrafo~\ref{sec:prg:testing} dà gli
stessi passi per QEMU e per la scheda.

\begin{shell}[caption={Compilare \fname{l05.S} ed eseguirlo su \prgqt{}.},label=lst:lab:l05:run]
# 1. start PRG32-QT on this computer (it listens on port 8080)

# 2. in the PRG32 folder, with ESP-IDF active: build the cartridge
python3 -m prg32 cartridge build ~/arc12/code/l05.S --portable \
    --entry-prefix l05 --name l05 --out l05.prg32

# 3. send it to PRG32-QT and start it
python3 -m prg32 esp32c6 upload-and-run l05.prg32 \
    --url http://localhost:8080 --slot cart0

# 4. look, press, compare with the previous slide
\end{shell}
